\subsection{Control de adopción y respuesta operacional}\label{sec:sd8-adopcion-respuesta}
Esta aplicación complementa AC-R-04/13/14/15 y las acciones de innovaciones, con referentes funcionales, ejecutores y suplencias de T-15. El titular funcional conserva decisiones sobre custodia, escuela, faena y acceso; Synaptix presta soporte técnico. Ante riesgo para personas o custodia se protege primero la actividad y se activa BC-03; el indicador no demora una emergencia ni permite ignorar una autorización faltante.

Proyecto y Calidad revisan en cada comité quincenal los pendientes operacionales y consolidan mensualmente adopción por perfil, turno, sitio y función. Se conserva población/versión, período, numerador, denominador, exclusiones justificadas y casos adversos. Cobertura de usuarios activos es usuarios habilitados que ejecutaron una tarea aplicable sobre usuarios habilitados con tarea programada; el login no cuenta como ejecución. Uso correcto es tareas con hecho válido completo sobre tareas aplicables, incluyendo fallidas y pendientes. Duplicación/retorno manual es hechos con doble captura o uso ordinario de mecanismo anterior sobre hechos aplicables; los folios de contingencia autorizada se muestran aparte con estado de incorporación, sin ocultar su carga ni descartarlos por dificultad. Satisfacción se reporta por respuestas y población convocada, con no respuesta visible y sin atribuir satisfacción a quien no respondió. \parencite[art.~89; p.~45]{sd8-ba}

Como criterio propuesto, ninguna tarea crítica se habilita sin usuario competente/autorizado y todo hecho crítico debe quedar íntegro y conciliado; un hecho perdido, duplicado, una entrega sin autorización o retorno ordinario a cuaderno genera acción inmediata. Para las funciones adicionales se compara la cobertura/uso con el período anterior y con la meta del piloto de SD13: caída, abandono o beneficio no evaluable obliga a contrastar causas y revisar el tratamiento, sin redefinir el universo para aparentar éxito. No se inventa una tasa de adopción observada. Cada corrección identifica ID, causa, perfil, pantalla/proceso, dueño, plazo anterior a la próxima tarea crítica y prueba de eficacia. Se comprueba fricción/accesibilidad, se corrige flujo, se imparte recuperación de formación con capacidad asignada y se reevalúa; capacitar de nuevo sin verificar la tarea no cierra el riesgo.

Proyecto acuerda validación y relevo antes de campaña/corte; ausencia del conocedor del legado o referente activa transferencia documentada y prueba por suplente, o bloqueo de la decisión sin evidencia. Fuera de campaña las altas mantienen evaluación propia. SRE prepara terminal/repuesto precargado y soporte en el canal y horario vigentes; el usuario no reinstala, parchea ni restaura. Un incidente de terreno conserva diagnóstico, aviso, trayecto, ventana y tiempo efectivo de sustitución; cuando falta acceso seguro o stock compatible se declara la brecha, sin afirmar reemplazo cumplido por enviar un repuesto.

Varadero/portería conservan plan, acreditación y formación antes de habilitar; Administración concilia hechos facturables sin duplicar; Escuela comprueba lista real y autorizaciones; Operaciones mantiene aviso, acuse y llamadas pendientes. Al recuperar se incorporan todos los folios y resultados tardíos con revisión de cantidades/IDs, datos protegidos y efectos, y se comprueba que ninguna persona/embarcación quedó indebidamente fuera del registro. Incumplimiento de capacidad o ventana se propaga a T-15/T-18 antes de reprogramar y consume una sola reserva del riesgo correspondiente. La cancelación global por marejada y la logística marina continúan sujetas a los controles abiertos de calendario y provisión; estas respuestas no acreditan por sí mismas su cabida.
